% Shared preamble for the HOF minor assignment (co-generation plant case)
\documentclass[10pt]{article}
\usepackage[a4paper,landscape,top=1.4cm,bottom=1.3cm,left=1.6cm,right=1.6cm,headheight=14pt,headsep=8pt,footskip=16pt]{geometry}
\usepackage[T1]{fontenc}
\usepackage[utf8]{inputenc}
\usepackage[scaled=0.95]{helvet}
\renewcommand{\familydefault}{\sfdefault}
\usepackage{microtype}
\usepackage{graphicx}
\usepackage[table]{xcolor}
\usepackage{multicol}
\usepackage{enumitem}
\usepackage{tabularx}
\usepackage{booktabs}
\usepackage{array}
\usepackage{fancyhdr}
\usepackage{titlesec}
\usepackage{textcomp}
\usepackage[hidelinks]{hyperref}

\definecolor{polblue}{RGB}{0,51,153}
\definecolor{redbox}{RGB}{249,211,211}
\definecolor{redline}{RGB}{192,0,0}
\definecolor{srcgray}{RGB}{110,110,110}
\definecolor{rowgray}{RGB}{242,242,242}

\setlength{\parindent}{0pt}
\setlength{\parskip}{3pt}
\setlength{\columnsep}{18pt}
\setlength{\columnseprule}{0.3pt}
\setlist{nosep,leftmargin=*,topsep=1pt}
\setlist[itemize]{label=\textbullet}

\titleformat{\section}{\color{polblue}\Large\bfseries}{}{0pt}{}
\titlespacing*{\section}{0pt}{0pt}{4pt}
\titleformat{\subsection}{\color{polblue}\normalsize\bfseries}{}{0pt}{}
\titlespacing*{\subsection}{0pt}{5pt}{2pt}

\pagestyle{fancy}
\fancyhf{}
\renewcommand{\headrulewidth}{0.4pt}
\fancyhead[L]{\small\color{polblue}Accident in a co-generation plant: HOF analysis}
\fancyhead[R]{\small\color{polblue}TRG Minor Assignment}
\fancyfoot[C]{\small\thepage}

% ---------- macros ----------
\newcommand{\ra}{\textrightarrow{}}
\newcommand{\src}[1]{{\color{srcgray}\scriptsize[#1]}}
\newcommand{\red}[1]{\colorbox{redbox}{\strut #1}}
\newcommand{\boxname}[1]{\textbf{\color{polblue}#1}}

% Template page: title, one-line event description, Mermaid diagram, key classification
\newcommand{\eventtemplate}[5]{%
  \clearpage
  \section{Event #1 \textbar{} #2}
  {\small #3}\par\vspace{2pt}
  \begin{center}
    \includegraphics[width=\linewidth,height=0.78\textheight,keepaspectratio]{diagrams/event#1.pdf}
  \end{center}
  \vfill
  {\footnotesize\setlength{\fboxsep}{4pt}\fcolorbox{polblue}{white}{\parbox{\dimexpr\linewidth-2\fboxsep-2\fboxrule}{#4}}}\par
  {\scriptsize\color{srcgray}#5}
}
% Justification page header
\newcommand{\justification}[1]{%
  \clearpage
  \section{Event #1 \textbar{} Justification of the template}
}
\newcommand{\legendline}{Red boxes: direct causal chain of the event (active error and/or failure in defence that produced it). White boxes: contributing conditions. Top box: latent factors influencing both pathways.}

% CREAM block: context, cognitive function, individual factors, error mode, antecedents
\newcommand{\cream}[6]{%
  \subsection{Cognitive view (CREAM, Hollnagel 1998)}
  \renewcommand{\arraystretch}{1.12}
  \begin{tabularx}{\linewidth}{>{\raggedright\arraybackslash\bfseries}p{2.9cm}X}
    \toprule
    Socio-technical context & #1\\
    Cognitive function & #2\\
    Individual factors & #3\\
    Error mode & #4\\
    General antecedents & #5\\
    Specific antecedent & #6\\
    \bottomrule
  \end{tabularx}
}
% Error classification block
\newcommand{\classification}[3]{%
  \subsection{Classification of the active error}
  \renewcommand{\arraystretch}{1.12}
  \begin{tabularx}{\linewidth}{>{\raggedright\arraybackslash\bfseries}p{2.9cm}X}
    \toprule
    Phenotype (Swain \& Guttman) & #1\\
    Reason's taxonomy & #2\\
    SRK level (Rasmussen) & #3\\
    \bottomrule
  \end{tabularx}
}
